%%%%%%%%%%%%%%%%%%%%%%%%%%%%%%%%%%%%%%%%%%%%%%%%%%%%%%%%%%%%%%%%%%%%%%
% Lớp 9 | Chương I. Phương trình và hệ hai phương trình bậc nhất hai ẩn | Giải bài toán bằng cách lập hệ phương trình
% Dạng 9D1?3-4: Toán thực tế (hóa học, vật lí, kinh tế)
%%%%%%%%%%%%%%%%%%%%%%%%%%%%%%%%%%%%%%%%%%%%%%%%%%%%%%%%%%%%%%%%%%%%%%


% nguồn: 2026-11-14_09_de_minh_hoa_tuyen_sinh_10_tphcm (nhập 14/11/2026)
\begin{bt}%[9D1V3-4]%[ma:11be934a]
% nguon: Đề thi tham khảo Tuyển sinh 10 TP.HCM – Phụ lục Sở GD&ĐT TP.HCM – Bài 6
Thép không gỉ Ferritic là họ thép hợp kim có chứa từ $12\%$ đến $27\%$ crôm. Một nhà máy luyện thép hiện có sẵn một lượng hợp kim thép chứa $10\%$ crôm và một lượng hợp kim thép chứa $30\%$ crôm. Giả sử trong quá trình luyện thép các nguyên liệu không bị hao hụt.
\begin{enumerate}[a)]
	\item Tính khối lượng hợp kim thép mỗi loại từ hai loại thép trên dùng để luyện được $500\text{ tấn}$ thép chứa $16\%$ crôm.
	\item Nhà máy dự định luyện ra loại thép không gỉ Ferritic từ $100\text{ tấn}$ thép chứa $10\%$ crôm và $x\text{ tấn}$ thép chứa $30\%$ crôm. Hỏi $x$ nằm trong khoảng nào?
\end{enumerate}
\loigiai{
\begin{enumerate}[a)]
	\item Gọi khối lượng hợp kim thép chứa $10\%$ crôm và chứa $30\%$ crôm cần dùng lần lượt là $u$ và $v$ (tấn, điều kiện: $0 < u, v < 500$).\\
	Theo đề bài, tổng khối lượng thép luyện được là $500\text{ tấn}$, ta có:
	$u + v = 500$.\\
	Lượng crôm nguyên chất có trong sản phẩm là $500 \cdot 16\% = 80\text{ (tấn)}$, ta có phương trình:
	$10\% u + 30\% v = 80 \Leftrightarrow 0{,}1u + 0{,}3v = 80 \Leftrightarrow u + 3v = 800$.\\
	Ta có hệ phương trình:
	$\begin{cases} u + v = 500 \\ u + 3v = 800 \end{cases} \Leftrightarrow \begin{cases} 2v = 300 \\ u = 500 - v \end{cases} \Leftrightarrow \begin{cases} v = 150 \\ u = 350 \end{cases}$ (thoả mãn điều kiện).\\
	Vậy cần dùng $350\text{ tấn}$ thép chứa $10\%$ crôm và $150\text{ tấn}$ thép chứa $30\%$ crôm.
	\item Khi trộn $100\text{ tấn}$ thép chứa $10\%$ crôm với $x\text{ tấn}$ thép chứa $30\%$ crôm ($x > 0$):
	\begin{itemize}
		\item Tổng khối lượng thép thu được là: $100 + x\text{ (tấn)}$.
		\item Tổng lượng crôm nguyên chất thu được là: $100\cdot 10\% + x\cdot 30\% = 10 + 0{,}3x\text{ (tấn)}$.
	\end{itemize}
	Tỉ lệ phần trăm crôm trong hợp kim mới là: $\dfrac{10 + 0{,}3x}{100 + x}$.\\
	Để hợp kim thu được là thép không gỉ Ferritic (chứa từ $12\%$ đến $27\%$ crôm), ta có bất phương trình kép:
	$12\% \le \dfrac{10 + 0{,}3x}{100 + x} \le 27\% \Leftrightarrow 0{,}12 \le \dfrac{10 + 0{,}3x}{100 + x} \le 0{,}27$.\\
	Vì $x > 0$ nên $100 + x > 0$, nhân cả hai vế với $(100 + x)$:
	\begin{itemize}
		\item $0{,}12(100 + x) \le 10 + 0{,}3x \Leftrightarrow 12 + 0{,}12x \le 10 + 0{,}3x \Leftrightarrow 0{,}18x \ge 2 \Leftrightarrow x \ge \dfrac{2}{0{,}18} = \dfrac{100}{9} \approx 11{,}11\text{ (tấn)}$.
		\item $10 + 0{,}3x \le 0{,}27(100 + x) \Leftrightarrow 10 + 0{,}3x \le 27 + 0{,}27x \Leftrightarrow 0{,}03x \le 17 \Leftrightarrow x \le \dfrac{17}{0{,}03} = \dfrac{1700}{3} \approx 566{,}67\text{ (tấn)}$.
	\end{itemize}
	Vậy khối lượng $x$ cần dùng nằm trong khoảng $\dfrac{100}{9} \le x \le \dfrac{1700}{3}$ (tấn), tức là khoảng từ $11{,}11\text{ tấn}$ đến $566{,}67\text{ tấn}$.
\end{enumerate}
}
\end{bt}

% nguồn: 2026-11-15_09_de_tham_khao_quan_1_de_1_va_de_2 (nhập 15/11/2026)
\begin{bt}%[9D1V3-4]%[ma:a859cf45]
% nguon: Đề tham khảo Tuyển sinh 10 TP.HCM 2025-2026 – Quận 1 – Đề 1 – Bài 6
Anh Bình đang có kế hoạch tiết kiệm để mua một chiếc xe ô tô cũ có giá $250\text{ triệu đồng}$ trong vòng một năm ($12\text{ tháng}$). Hiện tại, anh Bình đã có $50\text{ triệu đồng}$ trong tài khoản tiết kiệm.
\begin{enumerate}[a)]
	\item Hãy xác định số tiền tối thiểu anh Bình cần tiết kiệm mỗi tháng để đạt được mục tiêu mua xe sau $12\text{ tháng}$.
	\item Nếu vào tháng thứ $7$, anh Bình phải dùng $15\text{ triệu đồng}$ từ khoản tiết kiệm để gửi về hỗ trợ gia đình, thì trung bình mỗi tháng trong $12\text{ tháng}$ anh Bình cần tiết kiệm ít nhất bao nhiêu tiền để vẫn đạt được mục tiêu mua xe?
\end{enumerate}
\loigiai{
\begin{enumerate}[a)]
	\item Số tiền anh Bình còn thiếu để đủ $250\text{ triệu đồng}$ mua xe là:
	$250 - 50 = 200\text{ (triệu đồng)}$.\\
	Số tiền tối thiểu anh Bình cần tiết kiệm mỗi tháng trong $12\text{ tháng}$ là:
	$\dfrac{200}{12} = \dfrac{50}{3}\text{ (triệu đồng)} \approx 16{,}67\text{ triệu đồng} = 16\,666\,667\text{ đồng}$.
	\item Khi rút $15\text{ triệu đồng}$ gửi về cho gia đình, tổng số tiền anh Bình cần tiết kiệm trong cả năm là:
	$200 + 15 = 215\text{ (triệu đồng)}$.\\
	Số tiền ít nhất anh Bình cần tiết kiệm trung bình mỗi tháng là:
	$\dfrac{215}{12}\text{ (triệu đồng)} \approx 17{,}92\text{ triệu đồng} = 17\,916\,667\text{ đồng}$.
\end{enumerate}
}
\end{bt}

% nguồn: 2026-11-15_09_de_tham_khao_quan_1_de_1_va_de_2 (nhập 15/11/2026)
\begin{bt}%[9D1V3-4]%[ma:e7b4387c]
% nguon: Đề tham khảo Tuyển sinh 10 TP.HCM 2025-2026 – Quận 1 – Đề 2 – Bài 6
Chị Ba muốn mở một cửa hàng bán trà sữa với giá niêm yết là $24\,000\text{ đồng/ly}$. Nhân dịp khai trương, chị Ba cân nhắc hai hình thức khuyến mãi:
\begin{itemize}
	\item Phương án 1: Giảm giá $30\%$ cho mỗi ly trà sữa.
	\item Phương án 2: Mua $2$ ly được tặng thêm $1$ ly cùng loại.
\end{itemize}
Theo em, nếu hiệu ứng thu hút khách hàng của hai phương án là như nhau thì chị Ba nên chọn phương án nào để có lợi hơn cho cửa hàng? Vì sao?
\loigiai{
Xét số tiền thu được của cửa hàng đối với từng phương án trên mỗi đơn vị sản phẩm bán ra:
\begin{itemize}
	\item Theo Phương án 1 (giảm giá $30\%$ mỗi ly):\\
	Giá bán thực tế mỗi ly trà sữa là:
	$24\,000 \cdot (100\% - 30\%) = 24\,000 \cdot 0{,}7 = 16\,800\text{ (đồng/ly)}$.\\
	Do đó, mỗi ly bán ra cửa hàng thu về $16\,800\text{ đồng}$.
	\item Theo Phương án 2 (mua $2$ tặng $1$):\\
	Khi khách hàng nhận $3$ ly trà sữa thì chỉ phải trả tiền $2$ ly. Số tiền khách trả là:
	$2 \cdot 24\,000 = 48\,000\text{ (đồng)}$.\\
	Giá bán trung bình mỗi ly trà sữa mà cửa hàng giao cho khách là:
	$\dfrac{48\,000}{3} = 16\,000\text{ (đồng/ly)}$.\\
	Do đó, mỗi ly xuất ra cửa hàng chỉ thu về $16\,000\text{ đồng}$.
\end{itemize}
Vì $16\,800\text{ đồng} > 16\,000\text{ đồng}$ (chênh lệch $800\text{ đồng}$ trên mỗi ly) nên Phương án 1 mang lại doanh thu trên từng ly cao hơn so với Phương án 2.\\
Vậy chị Ba nên chọn Phương án 1 để có lợi hơn cho cửa hàng.
}
\end{bt}

% nguồn: 2026-11-15_09_de_tham_khao_quan_3_de_1_va_de_2 (nhập 15/11/2026)
\begin{bt}%[9D1V3-4]%[ma:20814d24]
% nguon: Đề tham khảo Tuyển sinh 10 TP.HCM 2025-2026 – Quận 3 – Đề 1 – Bài 6
Ba chiếc bình chứa nước có tổng dung tích là $105\text{ lít}$. Nếu đổ đầy nước vào bình thứ nhất rồi lấy toàn bộ lượng nước đó đem rót vào hai bình còn lại thì:
\begin{itemize}
	\item Hoặc làm đầy bình thứ hai và rót được thêm $\dfrac{4}{5}$ dung tích bình thứ ba.
	\item Hoặc làm đầy bình thứ ba và rót được thêm $\dfrac{5}{6}$ dung tích bình thứ hai.
\end{itemize}
Hãy xác định thể tích của mỗi bình (giả sử nước không bị hao hụt trong quá trình rót).
\loigiai{
Gọi dung tích của bình thứ nhất, bình thứ hai và bình thứ ba lần lượt là $x, y, z$ (lít; điều kiện: $x, y, z > 0$).\\
Theo đề bài, tổng dung tích ba bình là $105\text{ lít}$, nên ta có:
$x + y + z = 105$.\\
Từ hai trường hợp rót nước của bình thứ nhất, ta có hệ phương trình:
$\begin{cases} x = y + \dfrac{4}{5}z \\ x = z + \dfrac{5}{6}y. \end{cases}$\\
Suy ra:
$y + \dfrac{4}{5}z = z + \dfrac{5}{6}y \Leftrightarrow y - \dfrac{5}{6}y = z - \dfrac{4}{5}z \Leftrightarrow \dfrac{1}{6}y = \dfrac{1}{5}z \Leftrightarrow z = \dfrac{5}{6}y$.\\
Thay $z = \dfrac{5}{6}y$ vào biểu thức của $x$, ta được:
$x = y + \dfrac{4}{5}\left(\dfrac{5}{6}y\right) = y + \dfrac{2}{3}y = \dfrac{5}{3}y$.\\
Thay $x = \dfrac{5}{3}y$ và $z = \dfrac{5}{6}y$ vào phương trình tổng dung tích, ta có:
$\dfrac{5}{3}y + y + \dfrac{5}{6}y = 105 \Leftrightarrow \left(\dfrac{10 + 6 + 5}{6}\right)y = 105 \Leftrightarrow \dfrac{21}{6}y = 105 \Leftrightarrow \dfrac{7}{2}y = 105 \Leftrightarrow y = 30\text{ (lít)}$.\\
Từ đó suy ra:
\begin{itemize}
	\item $x = \dfrac{5}{3} \cdot 30 = 50\text{ (lít)}$.
	\item $z = \dfrac{5}{6} \cdot 30 = 25\text{ (lít)}$.
\end{itemize}
Vậy dung tích của bình thứ nhất là $50\text{ lít}$, bình thứ hai là $30\text{ lít}$ và bình thứ ba là $25\text{ lít}$.
}
\end{bt}

% nguồn: 2026-11-15_09_de_tham_khao_quan_4_de_1_va_de_2 (nhập 15/11/2026)
\begin{bt}%[9D1V3-4]%[ma:be4c260d]
% nguon: Đề tham khảo Tuyển sinh 10 TP.HCM 2025-2026 – Quận 4 – Đề 1 – Bài 6
Cô Mai nhập $210\text{ bao gạo}$ gồm hai loại: bao gạo loại 1 (mỗi bao nặng $25\text{ kg}$, bán với giá $22\,500\text{ đồng/kg}$) và bao gạo loại 2 (mỗi bao nặng $50\text{ kg}$, bán với giá $18\,200\text{ đồng/kg}$) với tổng chi phí nhập hàng là $99\,700\,000\text{ đồng}$. Do thời tiết ẩm ướt, $5\%$ khối lượng gạo loại 1 và $6\%$ khối lượng gạo loại 2 bị hỏng không thể bán được. Tổng khối lượng gạo còn lại bán được là $6{,}15\text{ tấn} = 6150\text{ kg}$. Hỏi sau khi bán hết toàn bộ số gạo còn lại, cô Mai lãi hay lỗ bao nhiêu tiền?
\loigiai{
Gọi số bao gạo loại 1 và loại 2 cô Mai đã nhập lần lượt là $x$ và $y$ ($x, y \in \mathbb{N}^*, x, y < 210$).\\
Vì tổng số bao gạo nhập về là $210$ bao nên ta có phương trình:
$x + y = 210$.\\
Tổng khối lượng gạo ban đầu của loại 1 là $25x\text{ (kg)}$, sau khi hỏng $5\%$ thì khối lượng bán được là:
$25x \cdot (100\% - 5\%) = 25x \cdot 0{,}95 = 23{,}75x\text{ (kg)}$.\\
Tổng khối lượng gạo ban đầu của loại 2 là $50y\text{ (kg)}$, sau khi hỏng $6\%$ thì khối lượng bán được là:
$50y \cdot (100\% - 6\%) = 50y \cdot 0{,}94 = 47y\text{ (kg)}$.\\
Vì tổng khối lượng gạo còn lại bán được là $6150\text{ kg}$ nên ta có phương trình:
$23{,}75x + 47y = 6150$.\\
Ta có hệ phương trình:
$\begin{cases} x + y = 210 \\ 23{,}75x + 47y = 6150 \end{cases} \Leftrightarrow \begin{cases} y = 210 - x \\ 23{,}75x + 47(210 - x) = 6150 \end{cases} \Leftrightarrow \begin{cases} 23{,}25x = 3720 \\ y = 210 - x \end{cases} \Leftrightarrow \begin{cases} x = 160 \\ y = 50 \end{cases}$ (thoả mãn điều kiện).\\
Khối lượng gạo loại 1 bán được là: $160 \cdot 23{,}75 = 3800\text{ (kg)}$.\\
Doanh thu từ gạo loại 1 là: $3800 \cdot 22\,500 = 85\,500\,000\text{ (đồng)}$.\\
Khối lượng gạo loại 2 bán được là: $50 \cdot 47 = 2350\text{ (kg)}$.\\
Doanh thu từ gạo loại 2 là: $2350 \cdot 18\,200 = 42\,770\,000\text{ (đồng)}$.\\
Tổng doanh thu thu được từ việc bán gạo là:
$85\,500\,000 + 42\,770\,000 = 128\,270\,000\text{ (đồng)}$.\\
Số tiền lãi của cô Mai sau khi trừ chi phí nhập hàng là:
$128\,270\,000 - 99\,700\,000 = 28\,570\,000\text{ (đồng)}$.\\
Vậy cô Mai có lãi và số tiền lãi là $28\,570\,000\text{ đồng}$.
}
\end{bt}

% nguồn: 2026-11-15_09_de_tham_khao_quan_4_de_1_va_de_2 (nhập 15/11/2026)
\begin{bt}%[9D1V3-4]%[ma:cd55e9a2]
% nguon: Đề tham khảo Tuyển sinh 10 TP.HCM 2025-2026 – Quận 4 – Đề 2 – Bài 6
Điểm trung bình môn Toán của $100$ học sinh trong hai lớp 9A và 9B là $7{,}608$. Hãy tính điểm trung bình của học sinh mỗi lớp, biết rằng số học sinh của lớp 9A nhiều hơn lớp 9B là $2$ học sinh và điểm trung bình của học sinh lớp 9B bằng $\dfrac{9}{10}$ điểm trung bình của học sinh lớp 9A.
\loigiai{
Gọi số học sinh của lớp 9A và lớp 9B lần lượt là $x$ và $y$ ($x, y \in \mathbb{N}^*, x, y < 100$).\\
Vì tổng số học sinh hai lớp là $100$ và lớp 9A nhiều hơn lớp 9B là $2$ học sinh, nên ta có hệ phương trình:
$\begin{cases} x + y = 100 \\ x - y = 2 \end{cases} \Leftrightarrow \begin{cases} x = 51 \\ y = 49. \end{cases}$\\
Gọi điểm trung bình của học sinh lớp 9A là $u$ ($u > 0$).\\
Khi đó điểm trung bình của học sinh lớp 9B là: $v = \dfrac{9}{10}u = 0{,}9u$.\\
Tổng điểm thi môn Toán của cả hai lớp là $100 \cdot 7{,}608 = 760{,}8\text{ (điểm)}$.\\
Ta có phương trình:
$51u + 49(0{,}9u) = 760{,}8 \Leftrightarrow 51u + 44{,}1u = 760{,}8 \Leftrightarrow 95{,}1u = 760{,}8 \Leftrightarrow u = \dfrac{760{,}8}{95{,}1} = 8$.\\
Suy ra điểm trung bình của lớp 9B là: $v = 0{,}9 \cdot 8 = 7{,}2$.\\
Vậy điểm trung bình của lớp 9A là $8{,}0$ điểm và lớp 9B là $7{,}2$ điểm.
}
\end{bt}

% nguồn: 2026-11-15_09_de_tham_khao_quan_5_de_1_va_de_2 (nhập 15/11/2026)
\begin{bt}%[9D1H3-4]%[ma:3e580092]
% nguon: Đề tham khảo Tuyển sinh 10 TP.HCM 2025-2026 – Quận 5 – Đề 1 – Bài 6
Nhân dịp tết Nguyên đán, một trường trung học cơ sở đã tổ chức cho học sinh các khối lớp 6, 7, 8 và 9 quyên góp tiền ủng hộ tự nguyện nhằm góp phần nhỏ giúp đỡ những học sinh có hoàn cảnh khó khăn. Kết thúc đợt quyên góp, bình quân mỗi học sinh của khối lớp 9 quyên góp được $22\,000$ đồng; bình quân mỗi học sinh của ba khối lớp 6, 7 và 8 quyên được $18\,000$ đồng. Hỏi bình quân mỗi học sinh của trường quyên góp được bao nhiêu tiền? Biết tổng số học sinh của ba khối lớp 6, 7 và 8 gấp 3 lần tổng số của học sinh khối lớp 9.
\loigiai{
Gọi $x$ là tổng số học sinh khối lớp 9 ($x \in \mathbb{N}^*$).\\
Vì tổng số học sinh của ba khối 6, 7 và 8 gấp 3 lần số học sinh khối 9 nên số học sinh ba khối 6, 7 và 8 là $3x$ (học sinh).\\
Tổng số học sinh toàn trường là: $x + 3x = 4x$ (học sinh).\\
Tổng số tiền học sinh khối lớp 9 quyên góp được là:
$22\,000 \cdot x$ (đồng).\\
Tổng số tiền học sinh ba khối 6, 7 và 8 quyên góp được là:
$18\,000 \cdot 3x = 54\,000x$ (đồng).\\
Tổng số tiền toàn trường quyên góp được là:
$22\,000x + 54\,000x = 76\,000x$ (đồng).\\
Bình quân mỗi học sinh toàn trường quyên góp được số tiền là:
$\dfrac{76\,000x}{4x} = 19\,000$ (đồng).\\
Vậy bình quân mỗi học sinh của trường quyên góp được $19\,000$ đồng.
}
\end{bt}

% nguồn: 2026-11-15_09_de_tham_khao_quan_5_de_1_va_de_2 (nhập 15/11/2026)
\begin{bt}%[9D1V3-4]%[ma:6385da8d]
% nguon: Đề tham khảo Tuyển sinh 10 TP.HCM 2025-2026 – Quận 5 – Đề 2 – Bài 4
Một xí nghiệp dự tính vận chuyển hàng bằng hai xe tải và đang phân vân giữa hai phương án là mua hẳn hai chiếc xe tải hoặc thuê hai xe tải (hai xe giống nhau). Nếu mua hai xe với giá mỗi xe là $200\,000\,000$ đồng thì mỗi ngày xí nghiệp phải trả tiền công cho tài xế và mua nhiên liệu cho một xe là $5\,000\,000$ đồng. Còn nếu thuê xe thì giá thuê một xe tải là $10\,000\,000$ đồng/ngày (đã bao gồm tiền công cho tài xế và nhiên liệu).
\begin{enumerate}[a)]
    \item Gọi $T$ là tổng số tiền xí nghiệp bỏ ra để vận chuyển hàng sau $n$ ngày. Lập công thức tính $T$ theo $n$ đối với mỗi phương án.
    \item Hỏi xí nghiệp vận chuyển hàng từ bao nhiêu ngày trở lên thì nên chọn phương án mua xe sẽ tiết kiệm hơn phương án thuê xe?
\end{enumerate}
\loigiai{
\begin{enumerate}[a)]
    \item Lập công thức tính tổng chi phí $T$ sau $n$ ngày vận chuyển:
    \begin{itemize}
        \item Phương án mua 2 xe tải:\\
        Chi phí mua 2 xe tải ban đầu là:
        $2 \cdot 200\,000\,000 = 400\,000\,000$ (đồng).\\
        Chi phí tài xế và nhiên liệu cho cả 2 xe mỗi ngày là:
        $2 \cdot 5\,000\,000 = 10\,000\,000$ (đồng/ngày).\\
        Tổng số tiền chi sau $n$ ngày theo phương án mua xe là:
        $T_1(n) = 400\,000\,000 + 10\,000\,000n$ (đồng).
        \item Phương án thuê 2 xe tải:\\
        Giá thuê cả 2 xe mỗi ngày là:
        $2 \cdot 10\,000\,000 = 20\,000\,000$ (đồng/ngày).\\
        Tổng số tiền chi sau $n$ ngày theo phương án thuê xe là:
        $T_2(n) = 20\,000\,000n$ (đồng).
    \end{itemize}
    \item Để phương án mua xe tiết kiệm hơn phương án thuê xe thì tổng chi phí mua xe phải nhỏ hơn tổng chi phí thuê xe:
    $T_1(n) < T_2(n) \Leftrightarrow 400\,000\,000 + 10\,000\,000n < 20\,000\,000n \Leftrightarrow 10\,000\,000n > 400\,000\,000 \Leftrightarrow n > 40$.\\
    Vì $n$ là số ngày nguyên dương nên $n \ge 41$.\\
    Vậy xí nghiệp vận chuyển hàng từ 41 ngày trở lên thì nên chọn phương án mua xe sẽ tiết kiệm hơn phương án thuê xe.
\end{enumerate}
}
\end{bt}

% nguồn: 2026-11-15_09_de_tham_khao_quan_5_de_1_va_de_2 (nhập 15/11/2026)
\begin{bt}%[9D1V3-4]%[ma:b1da4443]
% nguon: Đề tham khảo Tuyển sinh 10 TP.HCM 2025-2026 – Quận 5 – Đề 2 – Bài 6
Ông Năm dự định cải tạo một mảnh vườn hình chữ nhật có chiều dài bằng $2{,}5$ lần chiều rộng. Ông nhận thấy rằng, nếu đào một cái hồ có mặt hồ là hình chữ nhật thì sẽ chiếm mất $3\%$ diện tích mảnh vườn. Còn nếu giảm chiều dài mặt hồ $5\text{ m}$ và tăng chiều rộng mặt hồ $2\text{ m}$ thì mặt hồ là hình vuông và diện tích mặt hồ giảm được $20\text{ m}^2$. Tính chiều dài, chiều rộng của mảnh vườn hình chữ nhật.
\loigiai{
Gọi chiều rộng của mặt hồ ban đầu là $x\text{ (m)}$ và chiều dài của mặt hồ ban đầu là $y\text{ (m)}$ ($y > 5, x > 0, y > x$).\\
Khi giảm chiều dài $5\text{ m}$ và tăng chiều rộng $2\text{ m}$ thì mặt hồ trở thành hình vuông, ta có phương trình:
$y - 5 = x + 2 \Leftrightarrow y - x = 7 \Leftrightarrow y = x + 7$.\\
Diện tích mặt hồ ban đầu là $x \cdot y\text{ (m}^2\text{)}$.\\
Diện tích mặt hồ sau khi thay đổi kích thước là $(y - 5)(x + 2) = (x + 2)^2\text{ (m}^2\text{)}$.\\
Vì diện tích mặt hồ mới giảm $20\text{ m}^2$ so với ban đầu nên ta có phương trình:
$\begin{aligned}
xy - (x + 2)^2 = 20 &\Leftrightarrow x(x + 7) - (x^2 + 4x + 4) = 20 \\
&\Leftrightarrow x^2 + 7x - x^2 - 4x - 4 = 20 \\
&\Leftrightarrow 3x = 24 \Leftrightarrow x = 8\text{ (thỏa mãn)}.
\end{aligned}$\\
Suy ra chiều rộng của mặt hồ là $8\text{ m}$ và chiều dài của mặt hồ là $y = 8 + 7 = 15\text{ m}$.\\
Diện tích của mặt hồ là:
$S_{\text{hồ}} = 8 \cdot 15 = 120\text{ (m}^2\text{)}$.\\
Vì diện tích mặt hồ chiếm $3\%$ diện tích mảnh vườn nên diện tích mảnh vườn là:
$S_{\text{vườn}} = \dfrac{120}{3\%} = \dfrac{120}{0{,}03} = 4\,000\text{ (m}^2\text{)}$.\\
Gọi chiều rộng của mảnh vườn là $W\text{ (m)}$ ($W > 0$). Chiều dài của mảnh vườn là $2{,}5W\text{ (m)}$.\\
Ta có phương trình diện tích mảnh vườn:
$W \cdot 2{,}5W = 4\,000 \Leftrightarrow 2{,}5W^2 = 4\,000 \Leftrightarrow W^2 = 1\,600 \Rightarrow W = 40\text{ (m)}$.\\
Chiều dài của mảnh vườn là:
$2{,}5 \cdot 40 = 100\text{ (m)}$.\\
Vậy mảnh vườn có chiều rộng là $40\text{ m}$ và chiều dài là $100\text{ m}$.
}
\end{bt}

% nguồn: 2026-11-15_09_de_tham_khao_quan_5_de_3_va_quan_6_de_1 (nhập 15/11/2026)
\begin{bt}%[9D1H3-4]%[ma:ae28f26a]
% nguon: Đề tham khảo Tuyển sinh 10 TP.HCM 2025-2026 – Quận 5 – Đề 3 – Bài 6
Một bài thi trắc nghiệm gồm 20 câu hỏi, nếu trả lời đúng thì được 5 điểm, nếu trả lời sai thì bị trừ 2 điểm, nếu bỏ qua câu trả lời thì được 0 điểm.
\begin{enumerate}[a)]
    \item Bạn Dũng tham dự bài thi và được 47 điểm. Biết Dũng bỏ qua 5 câu. Hỏi Dũng trả lời đúng mấy câu và trả lời sai mấy câu?
    \item Bạn Hoa tham dự bài thi và được 59 điểm. Hỏi Hoa trả lời đúng mấy câu, trả lời sai mấy câu và bỏ qua bao nhiêu câu?
\end{enumerate}
\loigiai{
\begin{enumerate}[a)]
    \item Gọi $x, y$ lần lượt là số câu Dũng trả lời đúng và trả lời sai ($x, y \in \mathbb{N}$).\\
    Vì Dũng bỏ qua 5 câu nên tổng số câu Dũng đã làm là:
    $x + y = 20 - 5 = 15$.\\
    Số điểm Dũng đạt được là 47 điểm nên ta có phương trình:
    $5x - 2y = 47$.\\
    Ta có hệ phương trình:
    $\left\{\begin{aligned} &x + y = 15 \\ &5x - 2y = 47 \end{aligned}\right. \Leftrightarrow \left\{\begin{aligned} &2x + 2y = 30 \\ &5x - 2y = 47 \end{aligned}\right. \Leftrightarrow \left\{\begin{aligned} &7x = 77 \\ &y = 15 - x \end{aligned}\right. \Leftrightarrow \left\{\begin{aligned} &x = 11 \\ &y = 4. \end{aligned}\right.$\\
    Vậy bạn Dũng trả lời đúng 11 câu và trả lời sai 4 câu.
    \item Gọi $a, b, c$ lần lượt là số câu Hoa trả lời đúng, sai và bỏ qua ($a, b, c \in \mathbb{N}, a + b + c = 20$).\\
    Theo bài ra ta có phương trình điểm số:
    $5a - 2b = 59 \Leftrightarrow 5a = 59 + 2b$.\\
    Vì $b \ge 0$ nên $5a \ge 59 \Rightarrow a \ge 12$.\\
    Mà tổng số câu hỏi là 20 nên $a \le 20$.\\
    Mặt khác, $59 + 2b$ là số lẻ nên $5a$ phải là số lẻ, suy ra $a$ là số lẻ.\\
    Do đó $a \in \{13; 15; 17; 19\}$.
    \begin{itemize}
        \item Với $a = 13 \Rightarrow 2b = 5 \cdot 13 - 59 = 65 - 59 = 6 \Rightarrow b = 3$.\\
        Khi đó số câu bỏ qua là $c = 20 - (13 + 3) = 4$ (thỏa mãn).
        \item Với $a = 15 \Rightarrow 2b = 5 \cdot 15 - 59 = 75 - 59 = 16 \Rightarrow b = 8$.\\
        Khi đó $c = 20 - (15 + 8) = -3 < 0$ (loại).
        \item Với $a \ge 17$ thì $b \ge \dfrac{5 \cdot 17 - 59}{2} = 13 \Rightarrow a + b \ge 30 > 20$ (loại).
    \end{itemize}
    Vậy bạn Hoa trả lời đúng 13 câu, trả lời sai 3 câu và bỏ qua 4 câu.
\end{enumerate}
}
\end{bt}

% nguồn: 2026-11-15_09_de_tham_khao_quan_5_de_3_va_quan_6_de_1 (nhập 15/11/2026)
\begin{bt}%[9D1H3-4]%[ma:cc440386]
% nguon: Đề tham khảo Tuyển sinh 10 TP.HCM 2025-2026 – Quận 6 – Đề 1 – Bài 3
Cách Tết Trung thu 2 tháng, bà Hoa nhập 100 cái bánh dẻo hạt sen và 150 cái bánh chay mè đen hết $10\,230\,000$ đồng. Thấy bán được, bà Hoa nhập thêm 400 cái bánh dẻo hạt sen và 200 cái bánh chay mè đen (cùng giá nhập đợt đầu) hết $22\,920\,000$ đồng.
\begin{enumerate}[a)]
    \item Hỏi giá nhập của mỗi chiếc bánh dẻo hạt sen và bánh chay mè đen là bao nhiêu?
    \item Trước Tết Trung thu 3 ngày, bà Hoa đã bán hết bánh chay mè đen nhưng vẫn còn 100 cái bánh dẻo hạt sen. Để nhanh chóng bán hết số bánh đã nhập về, bà Hoa thực hiện chương trình khuyến mại như sau: mua cái bánh thứ nhất bằng giá niêm yết, cái bánh thứ hai giảm $30\%$ so với giá niêm yết, cái bánh thứ ba trở đi giảm $60\%$ so với giá niêm yết. Bạn Mai đến cửa hàng bà Hoa mua 5 cái bánh dẻo hạt sen phải trả $145\,000$ đồng. Hỏi giá niêm yết của mỗi cái bánh dẻo hạt sen là bao nhiêu?
\end{enumerate}
\loigiai{
\begin{enumerate}[a)]
    \item Gọi $x$ (đồng) là giá nhập của mỗi chiếc bánh dẻo hạt sen và $y$ (đồng) là giá nhập của mỗi chiếc bánh chay mè đen ($x, y > 0$).\\
    Theo đề bài, chi phí nhập bánh đợt đầu là:
    $100x + 150y = 10\,230\,000$.\\
    Chi phí nhập bánh đợt hai là:
    $400x + 200y = 22\,920\,000$.\\
    Ta có hệ phương trình:
    $\left\{\begin{aligned} &100x + 150y = 10\,230\,000 \\ &400x + 200y = 22\,920\,000 \end{aligned}\right. \Leftrightarrow \left\{\begin{aligned} &400x + 600y = 40\,920\,000 \\ &400x + 200y = 22\,920\,000 \end{aligned}\right. \Leftrightarrow \left\{\begin{aligned} &400y = 18\,000\,000 \\ &100x + 150y = 10\,230\,000 \end{aligned}\right. \Leftrightarrow \left\{\begin{aligned} &y = 45\,000 \\ &x = 34\,800. \end{aligned}\right.$\\
    Vậy giá nhập của mỗi chiếc bánh dẻo hạt sen là $34\,800$ đồng và bánh chay mè đen là $45\,000$ đồng.
    \item Gọi $p$ (đồng) là giá niêm yết của mỗi chiếc bánh dẻo hạt sen ($p > 0$).\\
    Khi bạn Mai mua 5 chiếc bánh dẻo hạt sen theo chương trình khuyến mại:
    \begin{itemize}
        \item Chiếc bánh thứ nhất có giá bằng $100\%$ giá niêm yết: $p$ (đồng).
        \item Chiếc bánh thứ hai được giảm $30\%$, có giá là: $(1 - 30\%)p = 0{,}7p$ (đồng).
        \item Ba chiếc bánh còn lại (từ chiếc thứ ba trở đi) được giảm $60\%$, mỗi chiếc có giá là: $(1 - 60\%)p = 0{,}4p$ (đồng).\\
        Số tiền trả cho 3 chiếc bánh này là: $3 \cdot 0{,}4p = 1{,}2p$ (đồng).
    \end{itemize}
    Tổng số tiền bạn Mai phải trả cho 5 chiếc bánh là:
    $p + 0{,}7p + 1{,}2p = 2{,}9p$ (đồng).\\
    Theo bài ra, số tiền Mai phải trả là $145\,000$ đồng nên ta có phương trình:
    $2{,}9p = 145\,000 \Leftrightarrow p = \dfrac{145\,000}{2{,}9} = 50\,000$ (đồng).\\
    Vậy giá niêm yết của mỗi chiếc bánh dẻo hạt sen là $50\,000$ đồng.
\end{enumerate}
}
\end{bt}

% nguồn: 2026-11-15_09_de_tham_khao_quan_6_de_2_va_quan_7_de_1 (nhập 15/11/2026)
\begin{bt}%[9D1H3-4]%[ma:f21f7a17]
% nguon: Đề tham khảo Tuyển sinh 10 TP.HCM 2025-2026 – Quận 6 – Đề 2 – Bài 3
Trên cánh đồng có diện tích $160\text{ ha}$ của một đơn vị sản xuất, người ta dành $60\text{ ha}$ để cấy thí điểm giống lúa mới, còn lại vẫn cấy giống lúa cũ. Khi thu hoạch, đầu tiên người ta gặt $8\text{ ha}$ giống lúa cũ và $7\text{ ha}$ giống lúa mới để đối chứng. Kết quả là $7\text{ ha}$ giống lúa mới cho thu hoạch nhiều hơn $8\text{ ha}$ giống lúa cũ là 2 tấn thóc. Biết rằng tổng số thóc (cả hai giống) thu hoạch cả vụ trên $160\text{ ha}$ là 860 tấn. Hỏi năng suất của mỗi giống lúa trên $1\text{ ha}$ là bao nhiêu tấn thóc?
\loigiai{
Diện tích cấy giống lúa cũ là:
$160 - 60 = 100\text{ (ha)}$.\\
Gọi $x$ (tấn/ha) là năng suất của giống lúa cũ và $y$ (tấn/ha) là năng suất của giống lúa mới ($x, y > 0$).\\
Theo đề bài, thu hoạch $7\text{ ha}$ giống lúa mới nhiều hơn $8\text{ ha}$ giống lúa cũ là 2 tấn thóc, ta có phương trình:
$7y - 8x = 2 \Leftrightarrow -8x + 7y = 2$.\\
Tổng sản lượng thóc thu hoạch trên cả cánh đồng $160\text{ ha}$ là 860 tấn, ta có phương trình:
$100x + 60y = 860 \Leftrightarrow 5x + 3y = 43$.\\
Ta có hệ phương trình:
$\left\{\begin{aligned} &-8x + 7y = 2 \\ &5x + 3y = 43 \end{aligned}\right. \Leftrightarrow \left\{\begin{aligned} &-24x + 21y = 6 \\ &35x + 21y = 301 \end{aligned}\right. \Leftrightarrow \left\{\begin{aligned} &59x = 295 \\ &3y = 43 - 5x \end{aligned}\right. \Leftrightarrow \left\{\begin{aligned} &x = 5 \\ &y = 6. \end{aligned}\right.$\\
Các giá trị $x = 5$ và $y = 6$ đều thỏa mãn điều kiện bài toán.\\
Vậy năng suất của giống lúa cũ là 5 tấn/ha và giống lúa mới là 6 tấn/ha.
}
\end{bt}

% nguồn: 2026-11-15_09_de_tham_khao_quan_6_de_2_va_quan_7_de_1 (nhập 15/11/2026)
\begin{bt}%[9D1H3-4]%[ma:3afb73c4]
% nguon: Đề tham khảo Tuyển sinh 10 TP.HCM 2025-2026 – Quận 7 – Đề 1 – Bài 6
Thép không gỉ Ferritic là họ thép hợp kim có chứa từ $12\%$ đến $27\%$ crôm. Một nhà máy luyện thép hiện có sẵn một lượng hợp kim thép chứa $10\%$ crôm và một lượng hợp kim thép chứa $30\%$ crôm (các nguyên liệu không bị hao hụt trong quá trình luyện).
\begin{enumerate}[a)]
    \item Tính khối lượng hợp kim thép mỗi loại từ hai loại thép trên dùng để luyện được $500$ tấn thép chứa $16\%$ crôm.
    \item Nhà máy dự định luyện ra loại thép không gỉ Ferritic từ $100$ tấn thép chứa $10\%$ crôm và $x$ tấn thép chứa $30\%$ crôm. Hỏi $x$ nằm trong khoảng nào?
\end{enumerate}
\loigiai{
\begin{enumerate}[a)]
    \item Gọi $x, y$ lần lượt là khối lượng hợp kim thép chứa $10\%$ crôm và chứa $30\%$ crôm cần dùng ($x, y > 0$, đơn vị tấn).\\
    Tổng khối lượng thép cần luyện là $500$ tấn, ta có phương trình:
    $x + y = 500$.\\
    Khối lượng crôm nguyên chất trong $500$ tấn thép chứa $16\%$ crôm là:
    $500 \cdot 16\% = 80\text{ (tấn)}$.\\
    Theo bài ra ta có phương trình:
    $10\% x + 30\% y = 80 \Leftrightarrow 0{,}1x + 0{,}3y = 80 \Leftrightarrow x + 3y = 800$.\\
    Ta có hệ phương trình:
    $\left\{\begin{aligned} &x + y = 500 \\ &x + 3y = 800 \end{aligned}\right. \Leftrightarrow \left\{\begin{aligned} &2y = 300 \\ &x = 500 - y \end{aligned}\right. \Leftrightarrow \left\{\begin{aligned} &y = 150 \\ &x = 350. \end{aligned}\right.$\\
    Vậy cần dùng $350$ tấn thép chứa $10\%$ crôm và $150$ tấn thép chứa $30\%$ crôm.
    \item Khi trộn $100$ tấn thép chứa $10\%$ crôm với $x$ tấn thép chứa $30\%$ crôm ($x > 0$):\\
    Tổng khối lượng hợp kim thu được là: $100 + x$ (tấn).\\
    Tổng khối lượng crôm nguyên chất là:
    $100 \cdot 10\% + x \cdot 30\% = 10 + 0{,}3x$ (tấn).\\
    Tỉ lệ phần trăm crôm trong hợp kim mới là:
    $\dfrac{10 + 0{,}3x}{100 + x}$.\\
    Để hợp kim mới là thép không gỉ Ferritic (chứa từ $12\%$ đến $27\%$ crôm), ta có bất phương trình kép:
    $12\% \le \dfrac{10 + 0{,}3x}{100 + x} \le 27\% \Leftrightarrow 0{,}12 \le \dfrac{10 + 0{,}3x}{100 + x} \le 0{,}27$.\\
    Vì $100 + x > 0$ nên:
    $\left\{\begin{aligned} &0{,}12(100 + x) \le 10 + 0{,}3x \\ &10 + 0{,}3x \le 0{,}27(100 + x) \end{aligned}\right. \Leftrightarrow \left\{\begin{aligned} &12 + 0{,}12x \le 10 + 0{,}3x \\ &10 + 0{,}3x \le 27 + 0{,}27x \end{aligned}\right. \Leftrightarrow \left\{\begin{aligned} &0{,}18x \ge 2 \\ &0{,}03x \le 17 \end{aligned}\right. \Leftrightarrow \left\{\begin{aligned} &x \ge \dfrac{100}{9} \approx 11{,}11 \\ &x \le \dfrac{1700}{3} \approx 566{,}67. \end{aligned}\right.$\\
    Vậy giá trị của $x$ nằm trong khoảng từ $\dfrac{100}{9}$ tấn đến $\dfrac{1700}{3}$ tấn (xấp xỉ từ $11{,}11$ tấn đến $566{,}67$ tấn).
\end{enumerate}
}
\end{bt}

% nguồn: 2026-11-15_09_de_tham_khao_quan_7_de_2_va_quan_8_de_1 (nhập 15/11/2026)
\begin{bt}%[9D1H3-4]%[ma:85ff838c]
% nguon: Đề tham khảo Tuyển sinh 10 TP.HCM 2025-2026 – Quận 7 – Đề 2 – Bài 3
Công ty A sản xuất những chiếc xe lăn cho người khuyết tật với số vốn đầu tư ban đầu là $500$ triệu đồng. Chi phí sản xuất ra một chiếc xe lăn là $2\,500\,000$ đồng và giá bán ra mỗi chiếc là $3\,000\,000$ đồng.
\begin{enumerate}[a)]
    \item Viết hàm số biểu diễn tổng số tiền đã đầu tư đến khi sản xuất ra được $x$ chiếc xe lăn và hàm số biểu diễn số tiền thu được khi bán ra $x$ chiếc xe lăn.
    \item Hỏi công ty cần phải bán bao nhiêu chiếc xe lăn mới có thể thu hồi được số vốn sản xuất ban đầu?
\end{enumerate}
\loigiai{
\begin{enumerate}[a)]
    \item Gọi $x$ là số chiếc xe lăn được sản xuất và bán ra ($x \in \mathbb{N}^*$).
    \begin{itemize}
        \item Tổng số tiền đã đầu tư khi sản xuất ra $x$ chiếc xe lăn gồm số vốn ban đầu và chi phí sản xuất:
        $C(x) = 500\,000\,000 + 2\,500\,000x\text{ (đồng)}$.
        \item Số tiền thu được khi bán ra $x$ chiếc xe lăn là:
        $R(x) = 3\,000\,000x\text{ (đồng)}$.
    \end{itemize}
    \item Để thu hồi được vốn đầu tư ban đầu thì số tiền thu được phải lớn hơn hoặc bằng tổng số tiền đã đầu tư:
    $R(x) \ge C(x) \Leftrightarrow 3\,000\,000x \ge 500\,000\,000 + 2\,500\,000x \Leftrightarrow 500\,000x \ge 500\,000\,000 \Leftrightarrow x \ge 1\,000$.\\
    Vậy công ty cần phải bán được ít nhất $1\,000$ chiếc xe lăn để thu hồi được vốn ban đầu.
\end{enumerate}
}
\end{bt}

% nguồn: 2026-11-15_09_de_tham_khao_quan_7_de_2_va_quan_8_de_1 (nhập 15/11/2026)
\begin{bt}%[9D1H3-4]%[ma:e445a02b]
% nguon: Đề tham khảo Tuyển sinh 10 TP.HCM 2025-2026 – Quận 7 – Đề 2 – Bài 5
Trong tháng 11, hai cửa hàng của một thương hiệu thời trang bán được tất cả $1\,100$ sản phẩm. Sang tháng 12, cửa hàng thứ nhất bán vượt mức $15\%$, cửa hàng thứ hai bán vượt mức $20\%$ so với tháng 11, do đó trong tháng 12 cả hai cửa hàng bán được $1\,295$ sản phẩm. Hỏi trong tháng 12 mỗi cửa hàng bán được bao nhiêu sản phẩm?
\loigiai{
Gọi $x, y$ lần lượt là số sản phẩm cửa hàng thứ nhất và cửa hàng thứ hai bán được trong tháng 11 ($x, y \in \mathbb{N}^*$, $x, y < 1\,100$).\\
Trong tháng 11 cả hai cửa hàng bán được $1\,100$ sản phẩm, ta có phương trình:
$x + y = 1\,100$.\\
Số sản phẩm bán thêm trong tháng 12 là:
$1\,295 - 1\,100 = 195\text{ (sản phẩm)}$.\\
Vì cửa hàng thứ nhất vượt mức $15\%$ và cửa hàng thứ hai vượt mức $20\%$ nên ta có phương trình:
$15\% x + 20\% y = 195 \Leftrightarrow 0{,}15x + 0{,}2y = 195 \Leftrightarrow 3x + 4y = 3\,900$.\\
Ta có hệ phương trình:
$\left\{\begin{aligned} &x + y = 1\,100 \\ &3x + 4y = 3\,900 \end{aligned}\right. \Leftrightarrow \left\{\begin{aligned} &4x + 4y = 4\,400 \\ &3x + 4y = 3\,900 \end{aligned}\right. \Leftrightarrow \left\{\begin{aligned} &x = 500 \\ &y = 1\,100 - 500 = 600. \end{aligned}\right.$\\
Các giá trị tìm được đều thỏa mãn điều kiện bài toán.\\
Trong tháng 12:
\begin{itemize}
    \item Cửa hàng thứ nhất bán được: $500 \cdot (1 + 15\%) = 500 \cdot 1{,}15 = 575$ (sản phẩm).
    \item Cửa hàng thứ hai bán được: $600 \cdot (1 + 20\%) = 600 \cdot 1{,}2 = 720$ (sản phẩm).
\end{itemize}
Vậy trong tháng 12, cửa hàng thứ nhất bán được $575$ sản phẩm và cửa hàng thứ hai bán được $720$ sản phẩm.
}
\end{bt}

% nguồn: 2026-11-15_09_de_tham_khao_quan_7_de_2_va_quan_8_de_1 (nhập 15/11/2026)
\begin{bt}%[9D1H3-4]%[ma:94d83e57]
% nguon: Đề tham khảo Tuyển sinh 10 TP.HCM 2025-2026 – Quận 8 – Đề 1 – Câu 6
Có ba thùng dầu đựng tổng cộng 123 lít dầu. Nếu đổ từ thùng thứ nhất sang thùng thứ hai 5 lít, rồi đổ từ thùng thứ hai sang thùng thứ ba 7 lít, tiếp tục đổ từ thùng thứ ba sang thùng thứ nhất 9 lít thì số dầu ở thùng thứ nhất sẽ ít hơn số dầu ở thùng thứ hai là 4 lít và bằng $\dfrac{2}{3}$ số dầu ở thùng thứ ba. Tính số lít dầu ở mỗi thùng lúc đầu.
\loigiai{
Gọi số lít dầu ban đầu ở thùng thứ nhất, thùng thứ hai và thùng thứ ba lần lượt là $x, y, z$ ($x, y, z > 0$ và $x + y + z = 123$).\\
Sau khi thực hiện các lần chuyển dầu:
\begin{itemize}
    \item Thùng thứ nhất bớt 5 lít và nhận 9 lít, số dầu còn lại là: $x - 5 + 9 = x + 4$ (lít).
    \item Thùng thứ hai nhận 5 lít và bớt 7 lít, số dầu còn lại là: $y + 5 - 7 = y - 2$ (lít).
    \item Thùng thứ ba nhận 7 lít và bớt 9 lít, số dầu còn lại là: $z + 7 - 9 = z - 2$ (lít).
\end{itemize}
Theo đề bài:
Số dầu thùng thứ nhất ít hơn thùng thứ hai 4 lít, ta có:
$(y - 2) - (x + 4) = 4 \Leftrightarrow y - x - 6 = 4 \Leftrightarrow y = x + 10$.\\
Số dầu thùng thứ nhất bằng $\dfrac{2}{3}$ số dầu thùng thứ ba, ta có:
$x + 4 = \dfrac{2}{3}(z - 2) \Leftrightarrow 3(x + 4) = 2z - 4 \Leftrightarrow 2z = 3x + 16 \Leftrightarrow z = 1{,}5x + 8$.\\
Thay các biểu thức của $y$ và $z$ theo $x$ vào phương trình tổng số dầu:
$x + (x + 10) + (1{,}5x + 8) = 123 \Leftrightarrow 3{,}5x + 18 = 123 \Leftrightarrow 3{,}5x = 105 \Leftrightarrow x = 30$.\\
Từ đó suy ra:
$y = 30 + 10 = 40\text{ (lít)}$,\\
$z = 1{,}5 \cdot 30 + 8 = 53\text{ (lít)}$.\\
Các giá trị tìm được đều thỏa mãn bài toán.\\
Vậy ban đầu thùng thứ nhất có 30 lít dầu, thùng thứ hai có 40 lít dầu và thùng thứ ba có 53 lít dầu.
}
\end{bt}
